\chapter{Rendre testable un logiciel legacy : résorber les régressions et s'en prémunir}
\label{chap:testabilite}
% Chapitre-chapeau du volet 2 : framework (t3), extension VS Code (t4), portage (t5).

Ce deuxième chapitre nous ramène chez Free Mobile et rassemble les trois
missions consacrées à la \textbf{testabilité} du \emph{Valo Engine} et de ses
forks. Il répond à la forme de dette dominante identifiée au
chapitre~\ref{chap:veille} : la \textbf{dette de test}, c'est-à-dire
l'absence de filet de sécurité automatisé autour d'un logiciel pourtant
critique pour la facturation : le \emph{Valo Engine} calcule le coût de
chaque appel, et une régression ou une panne sur ce maillon peut engendrer
des pertes \textbf{de l'ordre du million d'euros en une semaine}.

L'importance des tests dans un contexte legacy dépasse la seule détection
d'anomalies. Sans filet de non-régression, chaque modification (correctif,
évolution, et à plus forte raison refactoring) est un pari : la seule
validation possible est manuelle, donc partielle, coûteuse et non rejouable.
Les régressions se découvrent alors en production, là où leur coût est
maximal. Surtout, la dette de test verrouille toutes les autres : on ne peut
résorber une dette de code ou d'architecture qu'en modifiant le logiciel, ce
que l'on ne s'autorise pas sans filet. Rendre un logiciel legacy testable est
donc le \textbf{prérequis de tout remboursement} : c'est l'investissement qui
rend les autres possibles.

Encore faut-il que ce filet soit fiable et adopté. Les trois missions de ce
chapitre s'enchaînent dans cette logique : concevoir un framework de test
d'intégration et de non-régression à scénarios déclaratifs
(section~\ref{chap:t3}), lever les freins à son adoption par l'équipe grâce à
une extension VS~Code dédiée (section~\ref{chap:t4}), puis étendre ce filet
aux six forks legacy du moteur par un portage de l'architecture du framework
(section~\ref{chap:t5}).
